\chapter{Alcance del proyecto}\label{alcance}

\section{Objetivo}

Este proyecto está pensado para servir de iniciación e introducción a una de las propuestas del NIST, Super-Singular Key Exchange o SIKE. Este protocolo toma de punto de inicio las curvas elípticas supersingulares y adapta los protocolos de Diffie-Hellman y Elgamal con la intención de proporcionar la seguridad necesaria para afrontar los desafíos de la informática postcuántica.

Los objetivos de este trabajo son, pues:

\begin{itemize}
    \item Definir claramente todos aquellos conceptos matemáticos algo más avanzados que utiliza SIKE.
    \item Explicar el protocolo de manera sencilla para facilitar la comprensión de este.
    \item Proporcionar una interfaz que gráficamente muestre cómo funciona el protocolo.
\end{itemize}



\section{Funcionalidades}

El sitio contará con las siguientes funcionalidades:

\begin{itemize}
    \item Explicación paso a paso de cada protocolo.
    \item Enlaces de interés a algunos de los recursos más importantes.
\end{itemize}

\section{Diseño}

El proyecto estará diseñado únicamente para ordenadores debido a los objetivos del mismo y al gran número de partes móviles que este contiene. Para más información técnica, dirigirse al apartado \ref{tecnologias} Tecnologías.

\section{Plazos}

Al ser un Trabajo Fin de Grado del departamento de Matemáticas Aplicadas I, el plazo final de entrega será el 21 de junio de 2026.

\section{Entregables}

Los entregables finales del proyecto serán

\begin{itemize}
    \item \href{https://introduction-to-sike.vercel.app}{Sitio web activo y funcionando}.
    \item Memoria del proyecto (es decir, este documento).
    \item \href{https://github.com/albramvar1/introduction-to-sike}{Código fuente disponible}.
\end{itemize}